% ---------------------------------------------------------------------
%  08 · Conclusiones
% ---------------------------------------------------------------------
\seccion{Conclusiones}

% hallazgo: número, cifra, título, explicación
\newcommand{\hallazgo}[4]{%
  \begin{card}[height=17mm,valign=center]
    \begin{tabular}{@{}>{\centering\arraybackslash}m{9mm}@{\hspace{2mm}}>{\centering\arraybackslash}m{32mm}@{\hspace{3mm}}>{\raggedright\arraybackslash}m{0.56\linewidth}@{}}
      \badge{#1} & {\large\bfseries\boldmath\color{uisgreen}#2} &
      {\small\bfseries #3}\par\vspace{0.5mm}{\footnotesize\color{uisgray}#4}
    \end{tabular}
  \end{card}}

\begin{frame}{Tres hallazgos}
  \vspace{1mm}
  \hallazgo{1}{84 \% vs 55.8 \%}{El piso de potencia decide dónde vale la pena el DVFS}
    {En CPU casi toda la potencia es piso y no hay margen; en GPU sí}
  \hallazgo{2}{0.819}{Señales ligeras aproximan el régimen Roofline sin perfilar}
    {En familias no vistas, con inferencia de 97 µs; perfilar con Nsight alarga 60 veces la ejecución}
  \hallazgo{3}{$-2.8$ \%}{Clasificar por fase gana cuando hay cómputo sensible al reloj}
    {Frente a un reloj fijo en E-C; donde no hay esa fase (E-A), una regla fija rinde igual}
\end{frame}

% fila de respuesta: parte de la pregunta, GPU, CPU
\newcommand{\resp}[5][13]{%
  \fill[uiscard,rounded corners=1mm] (0,#2) rectangle (126,{#2+#1});
  \node[font=\small\bfseries,anchor=west,align=flush left,text width=39mm] at (3,{#2+#1/2}) {#3};
  \node[font=\small,anchor=west,align=flush left,text width=38mm] at (47,{#2+#1/2}) {#4};
  \node[font=\small,anchor=west,align=flush left,text width=36mm] at (88,{#2+#1/2}) {#5};}

\begin{frame}{Respuesta a la pregunta de investigación}
  \centering
  \begin{tikzpicture}[x=1mm,y=1mm]
    \node[font=\small\bfseries,text=uisgreen,anchor=west] at (49,53.5) {GPU};
    \node[font=\small\bfseries,text=uisred,anchor=west] at (90,53) {CPU};
    \resp[18]{32}{¿Reduce el\\consumo?}
      {\textbf{Sí:} $-3.8$ / $-3.9$ \% de EDP CPU+GPU\\[0.3mm]{\footnotesize\color{uisgray}solo GPU: $-8.0$ / $-6.3$ \%; con memoria larga}}
      {\textbf{No:} con un piso de 84 \% no hay margen}
    \resp{17}{¿La inferencia\\degrada el rendimiento?}
      {\textbf{No:} 97 µs por decisión; la duración no cambia}
      {\textbf{No es inferir} (16 µs): observar sube 5--9 \% la energía de CPU}
    \resp{2}{¿Frente a los\\gobernadores nativos?}
      {\textbf{Se mantiene} con \textit{performance} y \textit{powersave}}
      {Empeora el EDP con ambos}
    \fill[uisgreen,rounded corners=1mm] (0,-12) rectangle (126,-2);
    \node[font=\normalsize\bfseries,text=white] at (63,-7) {Recomendación: desplegar solo el proceso de GPU};
  \end{tikzpicture}
\end{frame}

% tarjeta de límite: título, texto
\newcommand{\limite}[2]{%
  \begin{card}[height=21mm,valign=top]
    {\small\bfseries\color{uisgreen}#1}\par\vspace{0.8mm}{\footnotesize #2}
  \end{card}}

\begin{frame}{Limitaciones: qué no afirma este trabajo}
  \vspace{1mm}
  \begin{columns}[T]
    \begin{column}{0.32\textwidth}
      \limite{Un solo nodo}{Xeon Gold 5315Y + A100. No se generaliza a otras CPU, GPU ni controladores}
      \limite{CPU acotada}{Solo la frecuencia de 6 núcleos delegados, sin turbo ni \textit{uncore}. No es la CPU completa}
    \end{column}
    \begin{column}{0.32\textwidth}
      \limite{Energía medida}{Paquetes de CPU y GPU. No incluye DRAM, placa, fuente ni ventilación: no es la energía del nodo}
      \limite{Gobernadores}{\textit{performance} y \textit{powersave}; \textit{ondemand} y \textit{schedutil} no existen con \texttt{intel\_pstate}}
    \end{column}
    \begin{column}{0.32\textwidth}
      \limite{Evidencia}{5 bloques: dirección consistente, no confirmación. E-C es una prueba de mecanismo construida}
      \limite{Modelo de GPU}{Entrenado por corrida y aplicado en ventanas de 3 s; con fases cortas o familias nuevas se abstiene}
    \end{column}
  \end{columns}
\end{frame}

% elemento de lista con viñeta verde
\newcommand{\vineta}[1]{\par\vspace{1.6mm}\noindent\textcolor{uisgreen}{\rule[0.5ex]{1.6mm}{1.6mm}}\hspace{2mm}#1}

\begin{frame}{Trabajo futuro}
  \vspace{4mm}
  \begin{cardhi}[width=0.86\textwidth,center]
    \small
    \vineta{Más familias externas \textit{memory-bound} en GPU, evaluadas una sola vez con el modelo congelado}
    \vineta{Repetir base, agente y F1 fijo en aplicaciones externas con cómputo sensible al reloj}
    \vineta{Reentrenar el modelo de GPU por ventana, con fronteras de fase de CUPTI}
    \vineta{Reducir primero el costo de observar en CPU, antes de revisar su política}
    \vineta{Probar en otras arquitecturas con un piso de potencia de CPU más bajo, donde el DVFS de CPU sí tenga margen}
    \vineta{Actuar también sobre la frecuencia del \textit{uncore} y de la memoria en CPU, no solo sobre la de los núcleos}
  \end{cardhi}
\end{frame}
